% WHAT A HYPHEN LIGATURES WITH.
%
% When a break falls between two letters that were set as one ligature,
% both halves are set again -- and THE HYPHEN IS ONE OF THE LETTERS THE
% FONT'S PROGRAM IS RUN OVER. trip.tfm ligatures `B' with `-' into `C', so
%
%   \hyphenchar\rip=`- \patterns{B1B}  on the word BB
%
% gives
%
%   ..\discretionary replacing 1
%   ...\rip C (ligature B-)
%   ..|\rip B
%   ..\rip A (ligature BB)
%
% -- ONE character before the break, not a `B' and a `-'. cmr10 ligatures no
% letter with a hyphen, which is why the words in the probe for
% breaking-inside-a-ligature never showed this.
%
% This engine set the letters again and then put a plain hyphen behind
% them. trip line 195 hyphenates a word of trip.tfm's own ligatures.
\catcode`\{=1 \catcode`\}=2
\tracingonline=1 \showboxdepth=99 \showboxbreadth=99
\font\rip=trip \rip \hyphenchar\rip=`-
\lccode`A=1 \lccode`B=2 \lccode`C=3 \uchyph=1
\patterns{A1B B1B B1C C1A}
\hsize=1000pt \parindent=0pt \pretolerance=-1 \tolerance=10000 \hbadness=10000
\message{[A a break inside a BB ligature]}
\setbox1=\vbox{\noindent A BB\par}
\showbox1
\message{[B one inside a CA ligature]}
\setbox1=\vbox{\noindent A CACA\par}
\showbox1
\message{[done]}
\end
